\documentclass{article}
\usepackage{amsmath}
\newcommand{\R}{\mathbb{R}}
\newcommand{\norm}[1]{\left\| #1 \right\|}
\newcommand{\greet}[2][Hello]{#1, #2!}
\DeclareMathOperator{\tr}{tr}
\newenvironment{note}{\textbf{Note:} }{}
\title{A \emph{small} test}
\begin{document}
\maketitle
\section{Introduction}
% a comment that disappears
This costs 5\% and uses \textit{italics}, \textbf{bold \emph{nested}} text. \greet{world} \greet[Ciao]{Sicilia}
Inline math $\norm{x} \in \R$ and \(a < b\).
\input{chapter1}
\begin{note}Be careful.\end{note}

\subsection*{Lists}
\begin{itemize}
  \item First
  \item Second with
  \begin{enumerate}
    \item inner one
    \item inner two
  \end{enumerate}
\end{itemize}

\begin{align}
  \tr(A) &= \sum_i a_{ii} \label{eq:tr} \\
  x &< y
\end{align}
See \eqref{eq:tr} and \cite[p.~3]{knuth} or \href{https://example.org}{a link}.
\begin{verbatim}
x = 100% {raw}
\end{verbatim}
\begin{theorem}Every \texttt{x} is fine.\end{theorem}
\end{document}
